\section{Comparison with METIS and Simulated Annealing}
\label{sec:comparison}

\subsection{Experimental Setup}

We compare three algorithms on three states that span the range of
congressional delegation sizes and geographic complexity:
\begin{itemize}
\item \textbf{NC} (North Carolina, $k=14$, $n \approx 2{,}660$ tracts):
  medium complexity, competitive two-party state with a mix of urban,
  suburban, and rural geography.
\item \textbf{WI} (Wisconsin, $k=8$, $n \approx 1{,}900$ tracts):
  smaller delegation, historically contested, elongated east-west shape.
\item \textbf{TX} (Texas, $k=38$, $n \approx 5{,}300$ tracts):
  large delegation, geographically diverse with major metropolitan areas
  and vast rural west.
\end{itemize}

All runs use 2020 Census data, geographic boundary-length edge weights
(the B.2 default~\citep{b2paper}), the standard bisection tree structure,
and a single seed ($s = s_0$, the content-derived seed).
The three algorithms compared are:

\begin{enumerate}
\item \textbf{METIS (baseline)}: Single \metis{} call per node, no refinement.
  This is the standard platform configuration~\citep{karypis1998}.
\item \textbf{SA(steps=10)}: \sa{} bisection with $\mathtt{steps\_per\_tract} = 10$,
  warm-started from the \metis{} solution~\citep{b19paper}.
\item \textbf{CVD(iters=20)}: \cvd{} bisection with
  $\mathtt{cvd\_iters} = 20$, k-farthest seed initialisation,
  approximate medoid (probe size 50), and 200-iteration rebalance.
\end{enumerate}

\paragraph{Metrics.}
For each (state, algorithm) pair we report:
\begin{itemize}
\item \textbf{Edge Cut (EC)}: Total weighted edge cut of the final $k$-district
  plan; lower is better.
\item \textbf{Polsby-Popper (PP)}: Mean PP compactness across all $k$ districts;
  higher is more compact.
  Baseline values from H.1~\citep{h1paper}: NC $\pp = 0.217$, WI $\pp = 0.198$,
  TX $\pp = 0.184$.
\item \textbf{Iterations}: For \cvd{}, the number of iterations to convergence.
\item \textbf{Runtime (s)}: Wall-clock time on a single core, 3.5 GHz CPU.
\end{itemize}

\subsection{Results}

Table~\ref{tab:main} presents the comparison.
\cvd{}(iters=20) achieves Polsby-Popper compactness competitive with
\metis{} --- within 5\% on NC and WI --- despite not optimising edge cuts.
The edge-cut values for \cvd{} are substantially higher than \metis{},
as expected: \cvd{} does not minimise edge cuts, and the Voronoi assignment
by BFS distance is not designed to follow geographic boundary lines.

\begin{table}[t]
\centering
\caption{%
  Algorithm comparison on NC ($k=14$), WI ($k=8$), TX ($k=38$).
  EC: total edge cut (lower is better).
  PP: mean Polsby-Popper compactness (higher is better).
  Iter: iterations to convergence (CVD only; ``---'' for others).
  Runtime: wall-clock seconds (single core).
  $\Delta$PP: PP relative to METIS baseline (percentage points).
}
\label{tab:main}
\renewcommand{\arraystretch}{1.25}
\begin{tabular}{@{}llrrrrr@{}}
\toprule
State & Algorithm & EC & PP & $\Delta$PP (\%) & Iter & Time (s) \\
\midrule
\multirow{3}{*}{NC ($k=14$)}
 & METIS (baseline)   & 4{,}821 & 0.217 & ---   & --- &  4.2 \\
 & SA (steps=10)      & 4{,}513 & 0.248 & +14.3 & --- & 51.7 \\
 & CVD (iters=20)     & 5{,}639 & 0.231 & +6.5  & 13  &  6.8 \\
\midrule
\multirow{3}{*}{WI ($k=8$)}
 & METIS (baseline)   & 2{,}103 & 0.198 & ---   & --- &  1.8 \\
 & SA (steps=10)      & 1{,}962 & 0.220 & +11.1 & --- & 20.3 \\
 & CVD (iters=20)     & 2{,}517 & 0.206 & +4.0  & 11  &  2.9 \\
\midrule
\multirow{3}{*}{TX ($k=38$)}
 & METIS (baseline)   & 11{,}440 & 0.184 & ---   & --- &  18.3 \\
 & SA (steps=10)      & 10{,}729 & 0.204 & +10.9 & --- & 227 \\
 & CVD (iters=20)     & 13{,}208 & 0.189 & +2.7  & 15  &  24.6 \\
\bottomrule
\end{tabular}
\end{table}

\subsection{Analysis of Results}

\paragraph{Edge cut.}
\cvd{} produces substantially higher edge cuts than \metis{} across all three
states (NC: $+17\%$, WI: $+20\%$, TX: $+15\%$).
This is expected: \cvd{} assigns tracts by geographic proximity, not by
following graph boundaries.
Voronoi boundaries frequently cross many adjacency edges because the bisection
boundary is a straight path through the graph (a geodesic in hop-count space)
rather than a minimum-cut boundary.

\paragraph{Polsby-Popper compactness.}
Despite higher edge cuts, \cvd{} achieves PP compactness \emph{above}
the METIS baseline on all three states.
The mechanism is geometric: Voronoi regions centred at well-separated seeds
tend to be roughly circular (or at least convex in graph space), which
correlates with geographic compactness.
METIS minimises edge cuts but produces plans shaped by the graph's
connectivity structure, which may follow non-compact county or road boundaries.
\cvd{}'s PP improvement over METIS is modest (2.7--6.5\% depending on state),
well below SA's 10.9--14.3\%, but \cvd{} achieves it without any
edge-cut-based objective.

\paragraph{Compactness ordering.}
The observed quality ordering is:
\[
  \text{METIS} \;<\; \text{CVD} \;<\; \text{SA(steps=10)},
\]
where $<$ denotes lower PP (worse compactness).
This ordering reflects the different compactness mechanisms: METIS minimises
boundary length (edge cut) but has no explicit compactness target; CVD
targets geometric proximity but not boundary minimisation; SA refines METIS
by escaping local optima in the edge-cut landscape.

\paragraph{Convergence.}
\cvd{} converges in 11--15 iterations across the three states, well within the
default $\mathtt{cvd\_iters} = 20$ budget.
NC converges in 13 iterations; WI (more elongated) in 11; TX (more complex
urban geography) in 15.
No convergence warnings were triggered.
Convergence is faster on elongated states (WI) because the k-farthest
initialisation places seeds near the eastern and western ends of the state,
requiring fewer iterations to stabilise.

\paragraph{Runtime.}
\cvd{}(iters=20) is approximately $1.6\times$ slower than \metis{} on NC
(6.8 s vs.\ 4.2 s) and approximately $7.6\times$ faster than SA(steps=10)
(6.8 s vs.\ 51.7 s).
The runtime ratio \cvd{}/\metis{} is approximately constant across states
(NC: 1.6$\times$, WI: 1.6$\times$, TX: 1.3$\times$), consistent with
the $O(n_{\mathrm{iter}} \times m)$ scaling predicted by
Proposition~\ref{prop:runtime} with $n_{\mathrm{iter}} \approx 13$ iterations.

\paragraph{Partisan outcomes.}
Table~\ref{tab:main} omits partisan seat counts because \cvd{} produces
geometrically different maps from METIS: the seat distributions vary at the
$\pm 1$ seat level on NC and WI, and at the $\pm 2$ seat level on TX.
This variability reflects the different geographic structure of \cvd{} plans
versus METIS plans, and is consistent with the B.7 finding that partisan
outcomes are sensitive to structure-layer choices.
A full 50-state partisan analysis of \cvd{} is deferred to G.14~\citep{g14paper}.

\subsection{The CVD Compactness Profile}

The Polsby-Popper results above suggest that \cvd{} occupies a distinct
position in the compactness-vs-boundary-cost trade-off space.
\metis{} minimises boundary cost (edge cut) and achieves modest compactness
as a side effect.
\sa{} further minimises boundary cost through annealing and achieves higher
compactness.
\cvd{} ignores boundary cost entirely and achieves compactness through
a completely different mechanism: geographic proximity.

This orthogonality is potentially useful in legal contexts.
A plaintiff challenging a map as an extreme outlier needs to characterise the
space of ``normal'' maps.
A space containing only edge-cut-minimising plans (METIS, SA) may not
capture the full range of geometrically reasonable alternatives.
Adding \cvd{} plans to the ensemble ensures that the ``natural partition''
criterion is represented, not just the ``boundary-minimising partition'' criterion.

\subsection{Relation to BisectionEnsemble}

\cvd{} is a Structure-layer algorithm; \be{}~\citep{h1paper} is a
SeedCompositor-layer strategy.
Running \be{}$(T=50)$ with \cvd{} bisectors would execute 50 independent
\cvd{} runs per node, each with a different base seed (and therefore a
different CVD node seed and different k-farthest initialisation), and select
the result with the lowest edge cut.
This combination is expected to improve \cvd{}'s PP further, since the
worst-EC \cvd{} plans tend to also be less compact.
This configuration is left for G.14~\citep{g14paper}.
